\documentclass[a4paper, 12pt]{article}

% 1. Подключаем преамбулу с пакетами и настройками
% Добавляем новые пакеты
\usepackage{fontspec}
\usepackage{polyglossia}

\setmainlanguage{russian}
\setotherlanguage{english}

\setmainfont{Times New Roman} 
\newfontfamily{\cyrillicfont}{Times New Roman}
\newfontfamily{\cyrillicfontsf}{Arial}
\newfontfamily{\cyrillicfonttt}{Courier New}

\usepackage{graphicx}         % Загружаем пакет для добавления различной графики
\usepackage{import}           % Загружаем пакет для продвинутого импорта
\usepackage{subcaption}       % Поддержка сетки из фигур
\usepackage{float}            % Улучшенное размещение изображений (опция H)
\usepackage{geometry}         % Управление отступами на странице
\usepackage{amsmath}          % Поддержка математики
\usepackage{amsfonts}         % Математические шрифты (\mathbb)
\usepackage{booktabs}         % Вертикальные разделители
\usepackage[dvipsnames]{xcolor}% Расширенная поддержка цветов
\usepackage[colorlinks]{hyperref}
\usepackage{indentfirst}      % Красная строка в первом абзаце
\usepackage{listings}         % Поддержка листингов кода
\usepackage[title,titletoc]{appendix}
\usepackage{title}         % Собственный титульный лист (vpdtitle.sty)

\setcounter{MaxMatrixCols}{20}

% Отступы на странице
\geometry{left=2cm, right=2cm, bottom=2cm, top=2cm}

% Настройка гиперссылок
\hypersetup{
    linkcolor=blue,
    filecolor=magenta,
    urlcolor=cyan
}

% Цвета для подсветки кода
\definecolor{codegreen}{rgb}{0,0.6,0}
\definecolor{codegray}{rgb}{0.5,0.5,0.5}
\definecolor{codepurple}{rgb}{0.58,0,0.82}
\definecolor{backcolour}{rgb}{0.95,0.95,0.92}

% Стиль листингов кода
\lstdefinestyle{codestyle}{
    backgroundcolor=\color{backcolour},
    commentstyle=\color{codegreen},
    keywordstyle=\color{magenta},
    numberstyle=\tiny\color{codegray},
    stringstyle=\color{codepurple},
    basicstyle=\ttfamily\footnotesize,
    breakatwhitespace=false,
    breaklines=true,
    captionpos=b,
    keepspaces=true,
    numbers=left,
    numbersep=5pt,
    showspaces=false,
    showstringspaces=false,
    showtabs=false,
    tabsize=2
}

\lstset{style=codestyle}
\lstset{extendedchars=\true}

% 2. Заполнение данных титульного листа
\title{Кодирование и шифрование}

% Заполните ваших авторов
\author{Ануфриев Андрей Сергеевич}
\affil{465029}
\author{Бахуров Николай Алексеевич}
\affil{504862}
\author{Раджапов Фёдор Сергеевич}
\affil{502544}

\supervisor{Зарина Ольга Андреевна}

\date{\today}

\begin{document}

% Генерация титульного листа
\maketitle

\newpage
\tableofcontents
\newpage

\section{Введение}
Целью данной лабораторной работы является практическое освоение методов линейной алгебры в задачах криптографии и теории кодирования. В работе рассматриваются линейные шифры (шифр Хилла), методы их криптоанализа и линейные корректирующие коды (код Хэмминга).

% --- ОСНОВНАЯ ЧАСТЬ ---
\section{Основная часть}

% Подключаем Задание 1
\subsection{Задание 1. Шифр Хилла}

\subsubsection{Алфавит и исходное сообщение}
Для выполнения задания был выбран алфавит из $n = 33$ символов (32 буквы русского алфавита + пробел). Полученные данные представлены в таблице \ref{tab:code_tab}.

\begin{table}[h!]
    \centering
    \caption{Таблица кодировки алфавита ($n = 33$)}
    \label{tab:code_tab}
    \begin{tabular}{|c|c||c|c||c|c||c|c|}
        \hline
        \textbf{Символ} & \textbf{Код} & \textbf{Символ} & \textbf{Код} & \textbf{Символ} & \textbf{Код} & \textbf{Символ} & \textbf{Код} \\ \hline
        А & 0  & И & 8  & Р & 16 & Ш & 24 \\ \hline
        Б & 1  & Й & 9  & С & 17 & Щ & 25 \\ \hline
        В & 2  & К & 10 & Т & 18 & Ъ & 26 \\ \hline
        Г & 3  & Л & 11 & У & 19 & Ы & 27 \\ \hline
        Д & 4  & М & 12 & Ф & 20 & Ь & 28 \\ \hline
        Е & 5  & Н & 13 & Х & 21 & Э & 29 \\ \hline
        Ж & 6  & О & 14 & Ц & 22 & Ю & 30 \\ \hline
        З & 7  & П & 15 & Ч & 23 & Я & 31 \\ \hline
        \textit{[пробел]} & 32 & & & & & & \\ \hline
    \end{tabular}
\end{table}

\noindent \textbf{Исходное сообщение (12 символов):} \texttt{СУИРНИКИ ТОП} \\
\textbf{Вектор численных номеров:} 
\[
P = \begin{pmatrix} 17 & 19 & 8 & 16 & 13 & 8 & 10 & 8 & 32 & 18 & 14 & 15 \end{pmatrix}
\]

\subsubsection{Матрицы-ключи и проверка их обратимости}
Были придуманны три матрицы-ключа размерами $2\times2$, $3\times3$ и $4\times4$. Чтобы матрица имела обратную в кольце вычетов $\mathbb{Z}_{33}$, необходимо, чтобы её определитель удовлетворял условию $\text{НОД}(\det(K), 33) = 1$.

\begin{enumerate}
    \item \textbf{Ключ $2\times2$ ($K_1$):}
    \[
    K_1 = \begin{pmatrix} 2 & 1 \\ 3 & 4 \end{pmatrix}
    \]
    Расчёт определителя: $\det(K_1) = 2 \cdot 4 - 1 \cdot 3 = 8 - 3 = 5$. \\
    Проверка условий: $\text{НОД}(5, 33) = 1$, значит обратный элемент $5^{-1} \pmod{33}$ существует. По расширенному алгоритму Евклида:
    \[
    33 = 5 \cdot 6 + 3, \quad 5 = 3 \cdot 1 + 2, \quad 3 = 2 \cdot 1 + 1
    \]
    Выражая единицу: $1 = 3 - 2 = 3 - (5 - 3) = 2 \cdot 3 - 5 = 2 \cdot (33 - 5 \cdot 6) - 5 = 2 \cdot 33 - 13 \cdot 5$. \\
    Следовательно, $-13 \equiv 20 \pmod{33}$. Таким образом, $5^{-1} \pmod{33} = 20$.

    \item \textbf{Ключ $3\times3$ ($K_2$):}
    \[
    K_2 = \begin{pmatrix} 1 & 2 & 1 \\ 0 & 1 & 3 \\ 2 & 1 & 1 \end{pmatrix}
    \]
    Расчёт определителя по правилу треугольников:
    \[
    \det(K_2) = 1(1 \cdot 1 - 3 \cdot 1) - 2(0 \cdot 1 - 3 \cdot 2) + 1(0 \cdot 1 - 1 \cdot 2) = -2 + 12 - 2 = 8
    \]
    По модулю 33: $\det(K_2) = 8 \pmod{33}$. $\text{НОД}(8, 33) = 1$. Обратный элемент: $8^{-1} \pmod{33} = 29$.

    \item \textbf{Ключ $4\times4$ ($K_3$):}
    \[
    K_3 = \begin{pmatrix} 1 & 0 & 0 & 2 \\ 0 & 1 & 0 & 0 \\ 0 & 0 & 2 & 1 \\ 0 & 0 & 1 & 1 \end{pmatrix}, \quad \det(K_3) = 1 \cdot (2 \cdot 1 - 1 \cdot 1) = 1 \pmod{33}
    \]
    Все определители по модулю 33 различны: $\det(K_1) \equiv 5$, $\det(K_2) \equiv 8$, $\det(K_3) \equiv 1$.
\end{enumerate}

\subsubsection{Шифрование}

Шифрование выполняется методом умножения вектор-строк на матрицу-ключ $C_i = (P_i \cdot K) \pmod{33}$.

\paragraph{А. Шифрование ключом $K_1$ ($2\times2$):}
\begin{enumerate}
    \item \textbf{Блок 1} $(17, 19)$ — буквы 'С', 'У':
    \[
    \begin{aligned}
        c_1 &= (17 \cdot 2 + 19 \cdot 3) = 34 + 57 = 91 \equiv 25 \pmod{33} \quad (\text{'Щ'}) \\
        c_2 &= (17 \cdot 1 + 19 \cdot 4) = 17 + 76 = 93 \equiv 27 \pmod{33} \quad (\text{'Ы'})
    \end{aligned}
    \]
    \item \textbf{Блок 2} $(8, 16)$ — буквы 'И', 'Р':
    \[
    \begin{aligned}
        c_3 &= (8 \cdot 2 + 16 \cdot 3) = 16 + 48 = 64 \equiv 31 \pmod{33} \quad (\text{'Я'}) \\
        c_4 &= (8 \cdot 1 + 16 \cdot 4) = 8 + 64 = 72 \equiv 6 \pmod{33} \quad (\text{'Ж'})
    \end{aligned}
    \]
    \item \textbf{Блок 3} $(13, 8)$ — буквы 'Н', 'И':
    \[
    \begin{aligned}
        c_5 &= (13 \cdot 2 + 8 \cdot 3) = 26 + 24 = 50 \equiv 17 \pmod{33} \quad (\text{'С'}) \\
        c_6 &= (13 \cdot 1 + 8 \cdot 4) = 13 + 32 = 45 \equiv 12 \pmod{33} \quad (\text{'М'})
    \end{aligned}
    \]
    \item \textbf{Блок 4} $(10, 8)$ — буквы 'К', 'И':
    \[
    \begin{aligned}
        c_7 &= (10 \cdot 2 + 8 \cdot 3) = 20 + 24 = 44 \equiv 11 \pmod{33} \quad (\text{'Л'}) \\
        c_8 &= (10 \cdot 1 + 8 \cdot 4) = 10 + 32 = 42 \equiv 9 \pmod{33} \quad (\text{'Й'})
    \end{aligned}
    \]
    \item \textbf{Блок 5} $(32, 18)$ — буквы '[пробел]', 'Т':
    \[
    \begin{aligned}
        c_9 &= (32 \cdot 2 + 18 \cdot 3) = 64 + 54 = 118 \equiv 19 \pmod{33} \quad (\text{'У'}) \\
        c_{10} &= (32 \cdot 1 + 18 \cdot 4) = 32 + 72 = 104 \equiv 5 \pmod{33} \quad (\text{'Е'})
    \end{aligned}
    \]
    \item \textbf{Блок 6} $(14, 15)$ — буквы 'О', 'П':
    \[
    \begin{aligned}
        c_{11} &= (14 \cdot 2 + 15 \cdot 3) = 28 + 45 = 73 \equiv 7 \pmod{33} \quad (\text{'З'}) \\
        c_{12} &= (14 \cdot 1 + 15 \cdot 4) = 14 + 60 = 74 \equiv 8 \pmod{33} \quad (\text{'И'})
    \end{aligned}
    \]
\end{enumerate}
\textbf{Зашифрованное сообщение 1:} \texttt{ЩЫЯЖСМЛЙУЕЗИ}

\paragraph{Б. Шифрование ключом $K_2$ ($3\times3$):}
\begin{enumerate}
    \item \textbf{Блок 1} $(17, 19, 8)$ — 'СУИ':
    \[
    \begin{aligned}
        c_1 &= (17 \cdot 1 + 19 \cdot 0 + 8 \cdot 2) = 17 + 0 + 16 = 33 \equiv 0 \pmod{33} \quad (\text{'А'}) \\
        c_2 &= (17 \cdot 2 + 19 \cdot 1 + 8 \cdot 1) = 34 + 19 + 8 = 61 \equiv 28 \pmod{33} \quad (\text{'Ь'}) \\
        c_3 &= (17 \cdot 1 + 19 \cdot 3 + 8 \cdot 1) = 17 + 57 + 8 = 82 \equiv 16 \pmod{33} \quad (\text{'Р'})
    \end{aligned}
    \]
    \item \textbf{Блок 2} $(16, 13, 8)$ — 'РНИ':
    \[
    \begin{aligned}
        c_4 &= (16 \cdot 1 + 13 \cdot 0 + 8 \cdot 2) = 16 + 16 = 32 \equiv 32 \pmod{33} \quad (\text{' '}) \\
        c_5 &= (16 \cdot 2 + 13 \cdot 1 + 8 \cdot 1) = 32 + 13 + 8 = 53 \equiv 20 \pmod{33} \quad (\text{'Х'}) \\
        c_6 &= (16 \cdot 1 + 13 \cdot 3 + 8 \cdot 1) = 16 + 39 + 8 = 63 \equiv 30 \pmod{33} \quad (\text{'Ю'})
    \end{aligned}
    \]
    \item \textbf{Блок 3} $(10, 8, 32)$ — 'КИ ':
    \[
    c_7 = 74 \equiv 8 \, (\text{'И'}), \quad c_8 = 60 \equiv 27 \, (\text{'Ы'}), \quad c_9 = 66 \equiv 0 \, (\text{'А'})
    \]
    \item \textbf{Блок 4} $(18, 14, 15)$ — 'ТОП':
    \[
    c_{10} = 48 \equiv 15 \, (\text{'П'}), \quad c_{11} = 65 \equiv 32 \, (\text{' '}), \quad c_{12} = 75 \equiv 9 \, (\text{'Й'})
    \]
\end{enumerate}
\textbf{Зашифрованное сообщение 2:} \texttt{АЬР ХЮИЫАП Й}

\paragraph{В. Шифрование ключом $K_3$ ($4\times4$):}
\begin{enumerate}
    \item \textbf{Блок 1} $(17, 19, 8, 16)$ — 'СУИР':
    \[
    \begin{aligned}
        c_1 &= 17 \cdot 1 + 19 \cdot 0 + 8 \cdot 0 + 16 \cdot 0 = 17 \equiv 17 \pmod{33} \quad (\text{'С'}) \\
        c_2 &= 17 \cdot 0 + 19 \cdot 1 + 8 \cdot 0 + 16 \cdot 0 = 19 \equiv 19 \pmod{33} \quad (\text{'У'}) \\
        c_3 &= 17 \cdot 0 + 19 \cdot 0 + 8 \cdot 2 + 16 \cdot 1 = 16 + 16 = 32 \equiv 32 \pmod{33} \quad (\text{' '}) \\
        c_4 &= 17 \cdot 2 + 19 \cdot 0 + 8 \cdot 1 + 16 \cdot 1 = 34 + 8 + 16 = 58 \equiv 25 \pmod{33} \quad (\text{'Щ'})
    \end{aligned}
    \]
    \item \textbf{Блок 2} $(13, 8, 10, 8)$ — 'НИКИ':
    \[
    c_5 = 29 \pmod{33} \, (\text{'Э'}), \quad c_6 = 8 \pmod{33} \, (\text{'И'}), \quad c_7 = 28 \pmod{33} \, (\text{'Ь'}), \quad c_8 = 42 \equiv 9 \pmod{33} \, (\text{'Й'})
    \]
    \item \textbf{Блок 3} $(32, 18, 14, 15)$ — ' ТОП':
    \[
    c_9 = 62 \equiv 29 \, (\text{'Э'}), \quad c_{10} = 18 \, (\text{'Т'}), \quad c_{11} = 43 \equiv 10 \, (\text{'К'}), \quad c_{12} = 93 \equiv 27 \, (\text{'Ы'})
    \]
\end{enumerate}
\textbf{Зашифрованное сообщение 3:} \texttt{СУ ЩЭИЬЙЭТКЫ}

\subsubsection{Расшифровка и дешифрация сообщений}
Для дешифрации сообщений находится обратная матрица $K^{-1} \pmod{33}$ через транспонированную матрицу алгебраических дополнений $A^T$:
\[
K^{-1} = \det(K)^{-1} \cdot A^T \pmod{33}
\]

\paragraph{А. Дешифрация ключом $K_1^{-1}$ ($2\times2$):}
1. $K_1 = \begin{pmatrix} 2 & 1 \\ 3 & 4 \end{pmatrix}$, $\det(K_1) = 5$, $5^{-1} \equiv 20 \pmod{33}$.
2. Алгебраические дополнения: $A_{11}=4, A_{12}=-3, A_{21}=-1, A_{22}=2$.
3. $A^T = \begin{pmatrix} 4 & -3 \\ -1 & 2 \end{pmatrix} \equiv \begin{pmatrix} 4 & 30 \\ 32 & 2 \end{pmatrix} \pmod{33}$.
4. Обратная матрица:
   \[
   K_1^{-1} = 20 \cdot \begin{pmatrix} 4 & 30 \\ 32 & 2 \end{pmatrix} = \begin{pmatrix} 80 & 600 \\ 640 & 40 \end{pmatrix} \equiv \begin{pmatrix} 14 & 6 \\ 13 & 7 \end{pmatrix} \pmod{33}
   \]
5. Расчёт дешифрации по блокам $P_i = (C_i \cdot K_1^{-1}) \pmod{33}$:
   \begin{itemize}
       \item $C_1 = (25, 27) \rightarrow (25\cdot14 + 27\cdot6, 25\cdot13 + 27\cdot7) = (512, 514) \equiv (17, 19) \pmod{33} \rightarrow \text{'С', 'У'}$
       \item $C_2 = (31, 6) \rightarrow (31\cdot14 + 6\cdot6, 31\cdot13 + 6\cdot7) = (470, 445) \equiv (8, 16) \pmod{33} \rightarrow \text{'И', 'Р'}$
       \item $C_3 = (17, 12) \rightarrow (17\cdot14 + 12\cdot6, 17\cdot13 + 12\cdot7) = (310, 305) \equiv (13, 8) \pmod{33} \rightarrow \text{'Н', 'И'}$
       \item $C_4 = (11, 9) \rightarrow (11\cdot14 + 9\cdot6, 11\cdot13 + 9\cdot7) = (208, 206) \equiv (10, 8) \pmod{33} \rightarrow \text{'К', 'И'}$
       \item $C_5 = (19, 5) \rightarrow (19\cdot14 + 5\cdot6, 19\cdot13 + 5\cdot7) = (296, 282) \equiv (32, 18) \pmod{33} \rightarrow \text{' ', 'Т'}$
       \item $C_6 = (7, 8) \rightarrow (7\cdot14 + 8\cdot6, 7\cdot13 + 8\cdot7) = (146, 147) \equiv (14, 15) \pmod{33} \rightarrow \text{'О', 'П'}$
   \end{itemize}
\textbf{Расшифрованное сообщение 1:} \texttt{СУИРНИКИ ТОП}

\paragraph{Б. Дешифрация ключом $K_2^{-1}$ ($3\times3$):}
1. $K_2 = \begin{pmatrix} 1 & 2 & 1 \\ 0 & 1 & 3 \\ 2 & 1 & 1 \end{pmatrix}$, $\det(K_2) = 8$, $8^{-1} \equiv 29 \pmod{33}$.
2. $A^T = \begin{pmatrix} -2 & -1 & 5 \\ 6 & -1 & -3 \\ -2 & 3 & 1 \end{pmatrix} \equiv \begin{pmatrix} 31 & 32 & 5 \\ 6 & 32 & 30 \\ 31 & 3 & 1 \end{pmatrix} \pmod{33}$.
3. Обратная матрица:
   \[
   K_2^{-1} = 29 \cdot \begin{pmatrix} 31 & 32 & 5 \\ 6 & 32 & 30 \\ 31 & 3 & 1 \end{pmatrix} \equiv \begin{pmatrix} 8 & 4 & 13 \\ 9 & 4 & 12 \\ 8 & 21 & 29 \end{pmatrix} \pmod{33}
   \]
4. Расчёт дешифрации по блокам $P_i = (C_i \cdot K_2^{-1}) \pmod{33}$:
   \begin{itemize}
       \item $C_1 = (0, 28, 16) \rightarrow (380, 448, 800) \equiv (17, 19, 8) \rightarrow \text{'С', 'У', 'И'}$
       \item $C_2 = (32, 20, 30) \rightarrow (676, 838, 1526) \equiv (16, 13, 8) \rightarrow \text{'Р', 'Н', 'И'}$
       \item $C_3 = (8, 27, 0) \rightarrow (307, 140, 428) \equiv (10, 8, 32) \rightarrow \text{'К', 'И', ' '}$
       \item $C_4 = (15, 32, 9) \rightarrow (480, 377, 840) \equiv (18, 14, 15) \rightarrow \text{'Т', 'О', 'П'}$
   \end{itemize}
\textbf{Расшифрованное сообщение 2:} \texttt{СУИРНИКИ ТОП}

\paragraph{В. Дешифрация ключом $K_3^{-1}$ ($4\times4$):}
1. $K_3 = \begin{pmatrix} 1 & 0 & 0 & 2 \\ 0 & 1 & 0 & 0 \\ 0 & 0 & 2 & 1 \\ 0 & 0 & 1 & 1 \end{pmatrix}$, $\det(K_3) = 1$, $1^{-1} \equiv 1 \pmod{33}$.
2. Обратная матрица $K_3^{-1} \equiv \begin{pmatrix} 1 & 0 & 0 & 31 \\ 0 & 1 & 0 & 0 \\ 0 & 0 & 1 & 32 \\ 0 & 0 & 32 & 2 \end{pmatrix} \pmod{33}$.
3. Расчёт дешифрации по блокам $P_i = (C_i \cdot K_3^{-1}) \pmod{33}$:
   \begin{itemize}
       \item $C_1 = (17, 19, 32, 25) \equiv (17, 19, 8, 16) \rightarrow \text{'С', 'У', 'И', 'Р'}$
       \item $C_2 = (29, 8, 28, 9) \equiv (13, 8, 10, 8) \rightarrow \text{'Н', 'И', 'К', 'И'}$
       \item $C_3 = (29, 18, 10, 27) \equiv (32, 18, 14, 15) \rightarrow \text{' ', 'Т', 'О', 'П'}$
   \end{itemize}
\textbf{Расшифрованное сообщение 3:} \texttt{СУИРНИКИ ТОП}

\subsubsection{Симуляция вредоносного вмешательства}
В каждое из зашифрованных сообщений было намеренно внесено по 3 случайные ошибки (замена символа из алфавита).

\paragraph{А. Сообщение 1 (Ключ $2\times2$):}
\begin{itemize}
    \item \textbf{Зашифрованное сообщение:} \texttt{ЩЫЯЖСМЛЙУЕЗИ}
    \item \textbf{Сообщение с вмешательством:} \texttt{\textbf{\underline{А}}ЫЯ\textbf{\underline{А}}СМЛЙ\textbf{\underline{А}}ЕЗИ} (заменены символы 1, 4 и 9)
    \item \textbf{Расчёт расшифровки по блокам:}
    \[
    \begin{aligned}
        C'_1 = (0, 27) \cdot K_1^{-1} &\equiv (30, 24) \pmod{33} \quad (\text{'Ю', 'Ш'}) \\
        C'_2 = (31, 0) \cdot K_1^{-1} &\equiv (5, 21) \pmod{33} \quad (\text{'Е', 'Х'}) \\
        C'_3 = (17, 12) \cdot K_1^{-1} &\equiv (13, 8) \pmod{33} \quad (\text{'Н', 'И'}) \\
        C'_4 = (11, 9) \cdot K_1^{-1} &\equiv (10, 8) \pmod{33} \quad (\text{'К', 'И'}) \\
        C'_5 = (0, 5) \cdot K_1^{-1} &\equiv (30, 2) \pmod{33} \quad (\text{'Ю', 'В'}) \\
        C'_6 = (7, 8) \cdot K_1^{-1} &\equiv (14, 15) \pmod{33} \quad (\text{'О', 'П'})
    \end{aligned}
    \]
    \item \textbf{Расшифрованное сообщение:} \texttt{ЮШЕХНИКИЮВОП}
\end{itemize}

\paragraph{Б. Сообщение 2 (Ключ $3\times3$):}
\begin{itemize}
    \item \textbf{Зашифрованное сообщение:} \texttt{АЬР ХЮИЫАП Й}
    \item \textbf{Сообщение с вмешательством:} \texttt{АЬР\textbf{\underline{ААА}}ЫАП Й} (заменены символы 4, 5 и 6 — полностью 2-й блок)
    \item \textbf{Расшифрованное сообщение:} \texttt{СУИАААКИ ТОП}
\end{itemize}

\paragraph{В. Сообщение 3 (Ключ $4\times4$):}
\begin{itemize}
    \item \textbf{Зашифрованное сообщение:} \texttt{СУ ЩЭИЬЙЭТКЫ}
    \item \textbf{Сообщение с вмешательством:} \texttt{СУ Щ\textbf{\underline{ААА}}ЙЭТКЫ} (заменены символы 5, 6 и 7)
    \item \textbf{Расшифрованное сообщение:} \texttt{СУИРНАШТ ТОП}
\end{itemize}

\paragraph{Симуляция вредоносного вмешательства} (по данным Python-вывода)
Заменяем 3 случайных символа (seed=42) в каждом зашифрованном сообщении.

\paragraph{А. Сообщение 1 (Ключ $2\times2$):}
\begin{itemize}
    \item \textbf{Оригинал:} \texttt{ЩЫЯЖСМЛЙУЕЗИ}
    \item \textbf{Изменено в позициях:} [0, 1, 10]
    \item \textbf{Искажённое:} \texttt{[О][П]ЯЖСМЛЙУЕ[С]И}
    \item \textbf{Дешифрованное (после повреждения):} \texttt{ЦЧИРНИКИ ТЦН}
    \item \textbf{Ожидаемое:} \texttt{СУИРНИКИ ТОП}
    \item \textbf{Ошибки:} поз.0: 'Ц' (22) вместо 'С' (17); поз.1: 'Ч' (23) вместо 'У' (19); поз.10: 'Ц' (22) вместо 'О' (14); поз.11: 'Н' (13) вместо 'П' (15)
\end{itemize}

\paragraph{Б. Сообщение 2 (Ключ $3\times3$):}
\begin{itemize}
    \item \textbf{Оригинал:} \texttt{АЬР ФЮИЫАП Й}
    \item \textbf{Изменено в позициях:} [1, 2, 8]
    \item \textbf{Искажённое:} \texttt{А[Ы][Е]   Ф Ю И Ы [В] П   Й}
    \item \textbf{Дешифрованное (после повреждения):} \texttt{УПЗРНИЪСШТОП}
    \item \textbf{Ожидаемое:} \texttt{СУИРНИКИ ТОП}
    \item \textbf{Ошибки:} поз.0: 'У' (19) вместо 'С' (17); поз.1: 'П' (15) вместо 'У' (19); поз.2: 'З' (7) вместо 'И' (8); поз.6: 'Ъ' (26) вместо 'К' (10); поз.7: 'С' (17) вместо 'И' (8); поз.8: 'Ш' (24) вместо ' ' (32)
\end{itemize}

\paragraph{В. Сообщение 3 (Ключ $4\times4$):}
\begin{itemize}
    \item \textbf{Оригинал:} \texttt{СУ ЩНИЬЛ ТКЫ}
    \item \textbf{Изменено в позициях:} [0, 1, 3]
    \item \textbf{Искажённое:} \texttt{[О][ ]   [Б] Н И Ь Л   Т К Ы}
    \item \textbf{Дешифрованное (после повреждения):} \texttt{О ЪННИКИ ТОП}
    \item \textbf{Ожидаемое:} \texttt{СУИРНИКИ ТОП}
    \item \textbf{Ошибки:} поз.0: 'О' (14) вместо 'С' (17); поз.1: ' ' (32) вместо 'У' (19); поз.2: 'Ъ' (26) вместо 'И' (8); поз.3: 'Н' (13) вместо 'Р' (16)
\end{itemize}

\paragraph{Локальный вывод по Заданию 1:}
В ходе выполнения задания был изучен алгоритм шифрования Хилла, основанный на матричном умножении и арифметике по модулю $n = 33$. По результатам проведенных вычислений сделаны следующие выводы:

1. \textbf{Условие расшифровки (обратимость ключа):} Чтобы сообщение можно было расшифровать, матрица-ключ должна быть обратимой по модулю 33. Для этого её определитель не должен иметь общих делителей с размером алфавита, то есть $\text{НОД}(\det(K), 33) = 1$. Если это условие не выполняется, обратную матрицу найти невозможно, и данные будут потеряны.

2. \textbf{Зависимость от размера ключа:} Увеличение размера матрицы (с $2\times2$ до $4\times4$) увеличивает длину шифруемого блока. При этом одна и та же буква исходного текста в разных местах сообщения шифруется разными символами, так как её значение зависит от соседних букв в блоке.

3. \textbf{Распространение ошибок при помехах:} Эксперимент с изменением символов показал, что замена хотя бы одной буквы в зашифрованном блоке приводит к искажению всех букв этого блока при расшифровке. Это происходит потому, что при умножении на обратную матрицу каждый символ исходного текста вычисляется с участием всех элементов зашифрованного блока.

\newpage

% Подключаем Задание 2
\subsection{Задание 2. Взлом шифра Хилла}

\subsubsection{Постановка задачи}
Смоделируем ситуацию, когда у нас есть два зашифрованных сообщения $C_1$ и $C_2$, полученных шифром Хилла с одним и тем же неизвестным ключом $K$ размера $2\times 2$, и известен открытый текст $P_1$, соответствующий $C_1$. Требуется найти $P_2$ — открытый текст второго сообщения.

Используется тот же алфавит из $n = 33$ символов (см. таблицу~\ref{tab:code_tab}), что и в Задании 1.

\subsubsection{Генерация ключа и шифрование}

Случайным образом сгенерируем ключ $K$ размера $2\times 2$ с обратимым по модулю 33 определителем:
\[
K = \begin{pmatrix} 5 & 3 \\ 2 & 7 \end{pmatrix}, \qquad \det(K) = 5\cdot 7 - 3\cdot 2 = 35 - 6 = 29.
\]
Проверка: $\text{НОД}(29, 33) = 1$ — ключ обратим по модулю 33.

\paragraph{Первое сообщение $P_1$ (12 символов):}
\texttt{КРИПТОГРАФИЯ}

Числовые коды (по таблице~\ref{tab:code_tab}):
\[
P_1 = (10,\; 16,\; 8,\; 15,\; 18,\; 14,\; 3,\; 16,\; 0,\; 19,\; 8,\; 30).
\]

\paragraph{Второе сообщение $P_2$ (12 символов):}
\texttt{СЕКРЕТНОЕСОО} (С,Е,К,Р,Е,Т,Н,О,Е,С,О,О).

Числовые коды:
\[
P_2 = (17,\; 5,\; 10,\; 16,\; 5,\; 18,\; 13,\; 14,\; 5,\; 17,\; 14,\; 14).
\]

\paragraph{Шифрование.} Каждые два символа образуют вектор-строку, умножаемую на $K$:
$C_i = (P_i \cdot K) \bmod 33$.

\noindent\textbf{Сообщение $P_1$:}
\[
\begin{aligned}
(10,16) \cdot K &= (10\cdot 5 + 16\cdot 2,\; 10\cdot 3 + 16\cdot 7) = (50+32,\; 30+112) = (82, 142) \equiv (16, 10) \pmod{33},\\
(8,15) \cdot K &= (8\cdot 5 + 15\cdot 2,\; 8\cdot 3 + 15\cdot 7) = (40+30,\; 24+105) = (70, 129) \equiv (4, 30) \pmod{33},\\
(18,14) \cdot K &= (18\cdot 5 + 14\cdot 2,\; 18\cdot 3 + 14\cdot 7) = (90+28,\; 54+98) = (118, 152) \equiv (19, 20) \pmod{33},\\
(3,16) \cdot K &= (3\cdot 5 + 16\cdot 2,\; 3\cdot 3 + 16\cdot 7) = (15+32,\; 9+112) = (47, 121) \equiv (14, 22) \pmod{33},\\
(0,19) \cdot K &= (0\cdot 5 + 19\cdot 2,\; 0\cdot 3 + 19\cdot 7) = (38, 133) \equiv (5, 1) \pmod{33},\\
(8,30) \cdot K &= (8\cdot 5 + 30\cdot 2,\; 8\cdot 3 + 30\cdot 7) = (40+60,\; 24+210) = (100, 234) \equiv (1, 3) \pmod{33}.
\end{aligned}
\]
Шифртекст $C_1$: коды $(16,10,4,30,19,20,14,22,5,1,1,3)$, что соответствует строке
\[
C_1 = \texttt{РКАЬУФОХЖББГ}.
\]

\noindent\textbf{Сообщение $P_2$:}
\[
\begin{aligned}
(17,5) \cdot K &= (17\cdot 5 + 5\cdot 2,\; 17\cdot 3 + 5\cdot 7) = (85+10,\; 51+35) = (95, 86) \equiv (29, 20) \pmod{33},\\
(10,16) \cdot K &= (82, 142) \equiv (16, 10) \pmod{33},\\
(5,18) \cdot K &= (5\cdot 5 + 18\cdot 2,\; 5\cdot 3 + 18\cdot 7) = (25+36,\; 15+126) = (61, 141) \equiv (28, 9) \pmod{33},\\
(13,14) \cdot K &= (13\cdot 5 + 14\cdot 2,\; 13\cdot 3 + 14\cdot 7) = (65+28,\; 39+98) = (93, 137) \equiv (27, 5) \pmod{33},\\
(5,17) \cdot K &= (5\cdot 5 + 17\cdot 2,\; 5\cdot 3 + 17\cdot 7) = (25+34,\; 15+119) = (59, 134) \equiv (26, 2) \pmod{33},\\
(14,14) \cdot K &= (14\cdot 5 + 14\cdot 2,\; 14\cdot 3 + 14\cdot 7) = (70+28,\; 42+98) = (98, 140) \equiv (32, 8) \pmod{33}.
\end{aligned}
\]
Шифртекст $C_2$: коды $(29,20,16,10,28,9,27,5,26,2,32,8)$, что соответствует
\[
C_2 = \texttt{ЭУРКЬЙЫЖЪВЯИ}.
\]

\subsubsection{Атака: восстановление ключа по известной паре $(P_1, C_1)$}

Пусть ключ имеет вид $K = \begin{pmatrix} a & b \\ c & d \end{pmatrix}$. Уравнение $P_1 \cdot K = C_1 \pmod{33}$ даёт 12 уравнений на 4 неизвестных. Возьмём первые два блока (первые 4 символа $P_1$ и $C_1$):

\[
\begin{cases}
10 a + 16 c \equiv 16 \pmod{33} \\
10 b + 16 d \equiv 10 \pmod{33} \\
8 a + 15 c \equiv 4 \pmod{33} \\
8 b + 15 d \equiv 30 \pmod{33}
\end{cases}
\]

Из первой пары (для $a, c$):
\[
\begin{cases}
10 a + 16 c \equiv 16 \\
8 a + 15 c \equiv 4
\end{cases}
\]

Решим через определитель:

Матрица системы для $(a, c)$:
\[
M = \begin{pmatrix} 10 & 16 \\ 8 & 15 \end{pmatrix}, \qquad \det(M) = 10\cdot 15 - 16\cdot 8 = 150 - 128 = 22.
\]
Проверка: $\text{НОД}(22, 33) = 11 \neq 1$ — матрица необратима! Значит, этот блок не подходит для прямого обращения.

Возьмём другой блок. Попробуем блоки 3–4 (символы 5–8): $P_1$: $(18,14)$ и $(3,16)$; $C_1$: $(19,20)$ и $(14,22)$.

\[
\begin{cases}
18 a + 14 c \equiv 19 \\
3 a + 16 c \equiv 14
\end{cases}
\]
\[
M = \begin{pmatrix} 18 & 14 \\ 3 & 16 \end{pmatrix}, \qquad \det(M) = 18\cdot 16 - 14\cdot 3 = 288 - 42 = 246 \equiv 246 - 7\cdot 33 = 246 - 231 = 15 \pmod{33}.
\]
$\text{НОД}(15, 33) = 3 \neq 1$ — снова необратима.

Возьмём блоки 1 и 5 (символы 1–2 и 9–10): $P_1$: $(10,16)$ и $(0,19)$; $C_1$: $(16,10)$ и $(5,1)$.

\[
\begin{cases}
10 a + 16 c \equiv 16 \\
0\cdot a + 19 c \equiv 5
\end{cases}
\]
Из второго: $19 c \equiv 5 \pmod{33}$. Найдём $19^{-1} \bmod 33$. $\text{НОД}(19,33)=1$.
$33 = 1\cdot 19 + 14$, $19 = 1\cdot 14 + 5$, $14 = 2\cdot 5 + 4$, $5 = 1\cdot 4 + 1$.
$1 = 5 - 4 = 5 - (14 - 2\cdot 5) = 3\cdot 5 - 14 = 3\cdot(19-14) - 14 = 3\cdot 19 - 4\cdot 14 = 3\cdot 19 - 4\cdot(33 - 19) = 7\cdot 19 - 4\cdot 33$.
Значит $19^{-1} \equiv 7 \pmod{33}$.
$c \equiv 7\cdot 5 = 35 \equiv 2 \pmod{33}$.

Подставим в первое: $10 a + 16\cdot 2 \equiv 16 \Rightarrow 10 a + 32 \equiv 16 \Rightarrow 10 a \equiv -16 \equiv 17 \pmod{33}$.
$10^{-1} \bmod 33$: $\text{НОД}(10,33)=1$. $33 = 3\cdot 10 + 3$, $10 = 3\cdot 3 + 1$, $1 = 10 - 3\cdot 3 = 10 - 3(33 - 3\cdot 10) = 10\cdot 10 - 3\cdot 33$. Значит $10^{-1} \equiv 10$.
$a \equiv 10\cdot 17 = 170 \equiv 170 - 5\cdot 33 = 170 - 165 = 5 \pmod{33}$.

Итак $a = 5$, $c = 2$. Аналогично для $(b,d)$ из уравнений:
\[
\begin{cases}
10 b + 16 d \equiv 10 \\
19 d \equiv 1
\end{cases}
\]
$d \equiv 7\cdot 1 = 7 \pmod{33}$.
$10 b + 16\cdot 7 \equiv 10 \Rightarrow 10 b + 112 \equiv 10 \Rightarrow 10 b \equiv 10 - 112 = -102 \equiv -102 + 4\cdot 33 = -102 + 132 = 30 \pmod{33}$.
$b \equiv 10\cdot 30 = 300 \equiv 300 - 9\cdot 33 = 300 - 297 = 3 \pmod{33}$.

\textbf{Найденный ключ:}
\[
K = \begin{pmatrix} 5 & 3 \\ 2 & 7 \end{pmatrix},
\]
что полностью совпадает с исходным ключом. Проверим на всех блоках $P_1$:

$(10,16)K \equiv (16,10)$, $(8,15)K \equiv (4,30)$, $(18,14)K \equiv (19,20)$ , $(3,16)K \equiv (14,22)$ , $(0,19)K \equiv (5,1)$ , $(8,30)K \equiv (1,3)$ .

\subsubsection{Нахождение обратной матрицы $K^{-1}$}

$\det(K) = 29$. Найдём $29^{-1} \bmod 33$.
$33 = 1\cdot 29 + 4$, $29 = 7\cdot 4 + 1$, $1 = 29 - 7\cdot 4 = 29 - 7(33 - 29) = 8\cdot 29 - 7\cdot 33$. Значит $29^{-1} \equiv 8 \pmod{33}$.

Алгебраические дополнения: $A_{11} = 7$, $A_{12} = -2$, $A_{21} = -3$, $A_{22} = 5$.
\[
A^T = \begin{pmatrix} 7 & -2 \\ -3 & 5 \end{pmatrix} \equiv \begin{pmatrix} 7 & 31 \\ 30 & 5 \end{pmatrix} \pmod{33}.
\]
\[
K^{-1} = 8 \cdot \begin{pmatrix} 7 & 31 \\ 30 & 5 \end{pmatrix} = \begin{pmatrix} 56 & 248 \\ 240 & 40 \end{pmatrix} \equiv \begin{pmatrix} 23 & 17 \\ 9 & 7 \end{pmatrix} \pmod{33}.
\]

Проверка:
\[
K \cdot K^{-1} = \begin{pmatrix} 5 & 3 \\ 2 & 7 \end{pmatrix} \begin{pmatrix} 23 & 9 \\ 17 & 7 \end{pmatrix} = \begin{pmatrix} 115+51 & 45+21 \\ 46+119 & 18+49 \end{pmatrix} = \begin{pmatrix} 166 & 66 \\ 165 & 67 \end{pmatrix} \equiv \begin{pmatrix} 1 & 0 \\ 0 & 1 \end{pmatrix} \pmod{33}.
\]

\subsubsection{Расшифровка второго сообщения $C_2$}

$C_2 = (29,20,16,10,28,9,27,5,26,2,32,8)$. Вычисляем $P_2 = C_2 \cdot K^{-1} \bmod 33$:

\[
\begin{aligned}
(29,20) \cdot K^{-1} &= (29\cdot 23 + 20\cdot 17,\; 29\cdot 9 + 20\cdot 7) = (667+340,\; 261+140) = (1007, 401) \\
&\equiv (1007 - 30\cdot 33,\; 401 - 12\cdot 33) = (1007-990,\; 401-396) = (17, 5) \pmod{33},\\
(16,10) \cdot K^{-1} &= (16\cdot 23 + 10\cdot 17,\; 16\cdot 9 + 10\cdot 7) = (368+170,\; 144+70) = (538, 214) \\
&\equiv (538 - 16\cdot 33,\; 214 - 6\cdot 33) = (538-528,\; 214-198) = (10, 16) \pmod{33},\\
(28,9) \cdot K^{-1} &= (28\cdot 23 + 9\cdot 17,\; 28\cdot 9 + 9\cdot 7) = (644+153,\; 252+63) = (797, 315) \\
&\equiv (797 - 24\cdot 33,\; 315 - 9\cdot 33) = (797-792,\; 315-297) = (5, 18) \pmod{33},\\
(27,5) \cdot K^{-1} &= (27\cdot 23 + 5\cdot 17,\; 27\cdot 9 + 5\cdot 7) = (621+85,\; 243+35) = (706, 278) \\
&\equiv (706 - 21\cdot 33,\; 278 - 8\cdot 33) = (706-693,\; 278-264) = (13, 14) \pmod{33},\\
(26,2) \cdot K^{-1} &= (26\cdot 23 + 2\cdot 17,\; 26\cdot 9 + 2\cdot 7) = (598+34,\; 234+14) = (632, 248) \\
&\equiv (632 - 19\cdot 33,\; 248 - 7\cdot 33) = (632-627,\; 248-231) = (5, 17) \pmod{33},\\
(32,8) \cdot K^{-1} &= (32\cdot 23 + 8\cdot 17,\; 32\cdot 9 + 8\cdot 7) = (736+136,\; 288+56) = (872, 344) \\
&\equiv (872 - 26\cdot 33,\; 344 - 10\cdot 33) = (872-858,\; 344-330) = (14, 14) \pmod{33}.
\end{aligned}
\]

Полученные коды: $(17,5,10,16,5,18,13,14,5,17,14,14)$, что соответствует слову
\[
P_2 = \texttt{СЕКРЕТНОЕСОО}.
\]

\textbf{Проверка первого сообщения:} $P_1 = C_1 \cdot K^{-1}$:
\[
(16,10)K^{-1} \equiv (10,16),\; (4,30)K^{-1} \equiv (8,15),\; (19,20)K^{-1} \equiv (18,14),\;
(14,22)K^{-1} \equiv (3,16),\;
(5,1)K^{-1} \equiv (0,19),\;
(1,3)K^{-1} \equiv (8,30).
\]
Восстановлено $P_1 = \texttt{КРИПТОГРАФИЯ}$.

\subsubsection{Вывод по Заданию 2}

Атака на основе известного открытого текста на шифр Хилла оказывается успешной, потому что шифр линеен. Зная одну пару (открытый текст, шифртекст), можно составить систему линейных уравнений относительно элементов ключевой матрицы и решить её по модулю $n$. Если какой-то блок матрицы $P_1$ оказывается необратимым по модулю 33, это не мешает атаке — достаточно взять другой блок или решить систему целиком.

\textbf{Основные выводы:}
\begin{enumerate}
    \item Шифр Хилла уязвим к атаке на основе известного открытого текста: одного сообщения из 12 символов достаточно для полного восстановления ключа $2\times 2$ (нужно всего 4 независимых уравнения).
    \item Для ключей большего размера потребуется больше символов известного открытого текста (для $m\times m$ — порядка $m^2$ символов).
    \item Необратимость отдельных блоков открытого текста по модулю 33 не является препятствием: можно использовать другие блоки или решать систему совместно.
    \item Увеличение размера ключа повышает стойкость, но принципиально не защищает от атаки на основе известного открытого текста.
\end{enumerate}

\newpage

% Подключаем Задание 3
\section{Задание 3. Код Хэмминга (7, 4)}

\subsection{Постановка задачи}
В данном задании требуется освоить корректирующий код Хэмминга с параметрами $(7, 4)$, который позволяет обнаруживать и исправлять одиночные ошибки в 7-битных кодовых словах при передаче 4 информационных бит. Необходимо составить матрицы генератора $G$ и проверочную $H$, закодировать интересное слово из 4 букв русского алфавита, сымитировать вредоносное вмешательство (замену 1, 2, 3 и 4 бит) и восстановить исходное сообщение.

\subsection{Алфавит и двоичное представление}
Для работы выбран русский алфавит из 32 букв (без буквы \texttt{Ё}). Каждой букве сопоставлен пятибитовый двоичный номер от \texttt{00000} до \texttt{11111}. Таблица соответствия приведена в таблице \ref{tab:alphabet}.

\begin{center}
\small
\begin{tabular}{|c|c|c|c|c|c|c|c|}
\hline
\textbf{А} & \textbf{Б} & \textbf{В} & \textbf{Г} & \textbf{Д} & \textbf{Е} & \textbf{Ж} & \textbf{З} \\
\hline
00000 & 00001 & 00010 & 00011 & 00100 & 00101 & 00110 & 00111 \\
\hline
\textbf{И} & \textbf{Й} & \textbf{К} & \textbf{Л} & \textbf{М} & \textbf{Н} & \textbf{О} & \textbf{П} \\
\hline
01000 & 01001 & 01010 & 01011 & 01100 & 01101 & 01110 & 01111 \\
\hline
\textbf{Р} & \textbf{С} & \textbf{Т} & \textbf{У} & \textbf{Ф} & \textbf{Х} & \textbf{Ц} & \textbf{Ч} \\
\hline
10000 & 10001 & 10010 & 10011 & 10100 & 10101 & 10110 & 10111 \\
\hline
\textbf{Ш} & \textbf{Щ} & \textbf{Ъ} & \textbf{Ы} & \textbf{Ь} & \textbf{Э} & \textbf{Ю} & \textbf{Я} \\
\hline
11000 & 11001 & 11010 & 11011 & 11100 & 11101 & 11110 & 11111 \\
\hline
\end{tabular}
\end{center}
\label{tab:alphabet}

Исходное слово из 4 букв: \textbf{ТЕСТ}. Каждая буква преобразуется в 5-битный код:

\begin{align}
    \text{Т} \rightarrow 10010, \\
    \text{Е} \rightarrow 00101, \\
    \text{С} \rightarrow 10001, \\
    \text{Т} \rightarrow 10010.
\end{align}

Полный двоичный код слова (20 символов): \texttt{10010001011000110010}.

\subsection{Матрицы генератора $G$ и проверочная $H$}
Для кода Хэмминга $(7, 4)$ выбраны стандартные матрицы. Матрица генератора $G$ размером $4 \times 7$ составляется как $G = [I_4 \;|\; P]$, где $P$ — матрица проверки чётности $4 \times 3$:

\[
G = \begin{pmatrix}
1 & 0 & 0 & 0 & 1 & 1 & 0 \\
0 & 1 & 0 & 0 & 1 & 0 & 1 \\
0 & 0 & 1 & 0 & 0 & 1 & 1 \\
0 & 0 & 0 & 1 & 1 & 1 & 1
\end{pmatrix}
\]

Проверочная матрица $H$ размером $3 \times 7$ составляется как $H = [P^T \;|\; I_3]$:

\[
H = \begin{pmatrix}
1 & 1 & 0 & 1 & 1 & 0 & 0 \\
1 & 0 & 1 & 1 & 0 & 1 & 0 \\
0 & 1 & 1 & 1 & 0 & 0 & 1
\end{pmatrix}
\]

Проверка корректности: $G \cdot H^T \equiv 0 \pmod{2}$, что подтверждает линейную зависимость строк и правильное построение кода.

\begin{center}
\fbox{\parbox{0.9\textwidth}{
\centering
\textbf{Матрица генератора $G$ (4$\times$7)} \\
$G = \begin{pmatrix}
1 & 0 & 0 & 0 & 1 & 1 & 0 \\
0 & 1 & 0 & 0 & 1 & 0 & 1 \\
0 & 0 & 1 & 0 & 0 & 1 & 1 \\
0 & 0 & 0 & 1 & 1 & 1 & 1
\end{pmatrix}$ \\
\vspace{4pt}
\textbf{Проверочная матрица $H$ (3$\times$7)} \\
$H = \begin{pmatrix}
1 & 1 & 0 & 1 & 1 & 0 & 0 \\
1 & 0 & 1 & 1 & 0 & 1 & 0 \\
0 & 1 & 1 & 1 & 0 & 0 & 1
\end{pmatrix}$
}}
\end{center}

\subsection{Процесс кодирования}
Полный 20-битный код разбивается на 5 блоков по 4 информационных бита. Каждый блок умножается на $G$ в поле $\mathbb{F}_2$ (по модулю 2), что даёт 7-битное кодовое слово. Результаты кодирования приведены в таблице \ref{tab:encode}.

\begin{table}[H]
    \centering
    \small
    \caption{Кодирование блоков с помощью матрицы $G$}
    \label{tab:encode}
    \begin{tabular}{|c|c|c|}
        \toprule
        Блок (4 бита) & Код (7 бит) & Десятичное значение \\
        \midrule
        1001 & 1001001 & 73 \\
        0001 & 0001111 & 15 \\
        0110 & 0110110 & 54 \\
        0011 & 0011100 & 28 \\
        0010 & 0010011 & 19 \\
        \bottomrule
    \end{tabular}
\end{table}

Полный закодированный код (35 бит) составляет:
\texttt{1001001 0001111 0110110 0011100 0010011}.

Процесс кодирования визуализирован на рисунке \ref{fig:flow}.

\begin{center}
\fbox{\parbox{0.95\textwidth}{
\small
\textbf{Схема кодирования} \quad ТЕСТ $\rightarrow$ 10010 00101 10001 10010 $\rightarrow$ 5 блоков по 4 бита \quad $\rightarrow$ \quad 1001001 \; 0001111 \; 0110110 \; 0011100 \; 0010011 \quad $\rightarrow$ \quad 35 бит
}}
\end{center}
\label{fig:flow}

\subsection{Симуляция вредоносного вмешательства}
Для демонстрации возможностей кода последовательно испорчены 1, 2, 3 и 4 бита в различных позициях 35-битного кода. Искажённые сообщения с выделенными изменёнными битами приведены в таблице \ref{tab:corrupt}.

\begin{table}[H]
    \centering
    \small
    \caption{Искажённые сообщения с выделенными ошибочными битами}
    \label{tab:corrupt}
    \begin{tabular}{|l|l|}
        \toprule
        Вмешательство & Искажённый код (выделены ошибки) \\
        \midrule
        1 бит (позиция 5) & \texttt{10010[1]1 0001111 0110110 0011100 0010011} \\
        2 бита (3, 15) & \texttt{100[0]001 0001111 011[0]110 0011100 00[1]0011} \\
        3 бита (8, 20, 30) & \texttt{1001001 0[1]01111 011011[1] 0011100 00[0]0011} \\
        4 бита (1, 10, 25, 32) & \texttt{1[1]01001 000[0]111 0110110 0011[0]00 0010[1]11} \\
        \bottomrule
    \end{tabular}
\end{table}

\subsection{Декодирование и исправление ошибок}
Декодирование выполняется с помощью матрицы $H$. Для каждого 7-битного блока вычисляется синдром $s = H \cdot c^T \pmod{2}$. Синдром однозначно указывает позицию ошибочного бита (если он один в блоке). Таблица соответствия синдромов приведена на рисунке \ref{fig:syndrome}.

\begin{center}
\fbox{\parbox{0.92\textwidth}{
\small
\textbf{Процесс исправления ошибок с матрицей $H$} \\
1. Вычисляем синдром: $s = H \cdot c^T \pmod{2}$ \\
2. Синдром указывает позицию ошибочного бита \\
3. Инвертируем найденный бит \\
4. Извлекаем информационные биты (первые 4) \\
\vspace{4pt}
\textbf{Таблица синдромов} \\
\begin{tabular}{|c|c|}
\hline
\textbf{Синдром} & \textbf{Позиция ошибки} \\
\hline
000 & нет ошибки \\
110 & бит 1 \\
101 & бит 2 \\
011 & бит 3 \\
111 & бит 4 \\
100 & бит 5 (p1) \\
010 & бит 6 (p2) \\
001 & бит 7 (p3) \\
\hline
\end{tabular}
}}
\end{center}
\label{fig:syndrome}

Для каждого случая вмешательства выполнена декодировка. Результаты приведены в таблице \ref{tab:decode}. При 1--3 ошибках в разных блоках код успешно исправляет все ошибки и восстанавливает исходное слово \textbf{ТЕСТ}. При 4 ошибках, распределённых по разным блокам, также достигается полное восстановление, что демонстрирует высокую устойчивость кода при правильном распределении ошибок.

\begin{table}[H]
    \centering
    \caption{Результаты декодирования после исправления ошибок}
    \label{tab:decode}
    \begin{tabular}{|l|l|l|}
        \toprule
        Вмешательство & Декодированное слово & Верно? \\
        \midrule
        1 бит & ТЕСТ & ДА \\
        2 бита & ТЕСТ & ДА \\
        3 бита & ТЕСТ & ДА \\
        4 бита & ТЕСТ & ДА \\
        \bottomrule
    \end{tabular}
\end{table}

\subsection{Расчёты исправления ошибок}
Приведём пример расчёта для блока с одной ошибкой. Пусть кодовый блок $c = [1, 0, 1, 1, 0, 0, 1]$ (ошибка в бите 2). Вычисляем синдром:

\[
s = H \cdot c^T = \begin{pmatrix}
1 & 1 & 0 & 1 & 1 & 0 & 0 \\
1 & 0 & 1 & 1 & 0 & 1 & 0 \\
0 & 1 & 1 & 1 & 0 & 0 & 1
\end{pmatrix}
\begin{pmatrix}
1 \\ 0 \\ 1 \\ 1 \\ 0 \\ 0 \\ 1
\end{pmatrix}
= \begin{pmatrix}
1 \oplus 0 \oplus 1 \oplus 0 \oplus 0 \\ 1 \oplus 1 \oplus 1 \oplus 0 \oplus 0 \\ 0 \oplus 0 \oplus 1 \oplus 0 \oplus 1
\end{pmatrix}
= \begin{pmatrix}
1 \\ 0 \\ 1
\end{pmatrix} = (101)_2.
\]

Синдром $(101)_2$ соответствует ошибке в бите 2. После инвертирования второго бита и извлечения первых 4 информационных бит получаем исходные данные блока.

\subsection{Выводы по заданию 3}
В ходе выполнения задания были составлены матрицы $G$ и $H$ для кода Хэмминга $(7, 4)$, выполнено кодирование и декодирование 20-битного сообщения, сымитировано вредоносное вмешательство с заменой 1--4 бит. Показано, что код позволяет обнаруживать и исправлять одиночные ошибки в каждом 7-битном блоке. При распределении ошибок по разным блокам код успешно восстанавливает исходное слово даже при 4 ошибках. Это демонстрирует эффективность корректирующих кодов для защиты информации при передаче.


\end{document}